% Section 2 — Research methodology
\section{Research Methodology}
\label{sec:method}

\subsection{End-to-end research workflow}
\label{sec:workflow}
The study runs as a notebook ladder: a research map (00) fixes the question;
environment control (01) freezes seeds and software; acquisition and audit
notebooks (02a, 02b) verify both datasets; exploratory analysis (03)
characterizes demand before any model is chosen; experimental design (04) locks
the protocol; a teaching notebook (05a) establishes time-series vocabulary;
then the model rungs execute in order --- baselines (05b), exponential
smoothing (06), Croston specialists (06c), ARIMA/SARIMA (07), LSTM (08). A
post-hoc hardening pass corrected nine classes of issues (leap-day calendar
gap, LSTM validation leakage, subset-comparability, metric methodology, shared
policy) and froze the evidence; see \cref{sec:hardening}. Every notebook maps
to report sections in the traceability matrix
(\texttt{traceability/notebook\_traceability.csv}).

\subsection{Data acquisition and audit}
\label{sec:acquisition}
M5 source files (\texttt{sales\_train\_evaluation.csv}, \texttt{calendar.csv})
were audited for dimensions (30,490 item--store series, day columns d\_1 onward
with a verified d-index alignment), and the Store file (\texttt{train.csv}) was
verified as a complete panel: every store--item--date combination present, no
duplicates, a single zero-sales entry. Both datasets were reindexed onto one
common calendar so that cross-dataset comparison never confounds ``model
effect'' with ``different history lengths.''

\subsection{Temporal experimental protocol}
\label{sec:protocol}
\Cref{tab:protocol} is the frozen design. Time-based splits only --- no
shuffling, ever. Training covers 2013-01-01 to 2015-10-31 (1,034 days);
validation 2015-11-01 to 2016-02-29 (121 days, leap day included after the
hardening correction); testing 2016-03-01 to 2016-05-22 (83 days). Forecasts
span $h=28$ days from each of 8 weekly rolling origins under an expanding
window: every model sees strictly the history before each origin, asserted in
code. All selection (windows, smoothing weights, ARIMA orders, Croston
alphas, LSTM stopping) uses validation data only; the test window is touched
once for final scoring. Global seed 42 governs all stochastic steps.

\begin{table}[H]
\centering
\caption{Frozen experimental protocol (source: \texttt{05\_experiments/config.json}).}
\label{tab:protocol}
\begin{tabular}{ll}
\toprule
Element & Frozen value \\
\midrule
Common window & 2013-01-01 to 2016-05-22 (1,238 days) \\
Train / validation / test & 1,034 / 121 / 83 days (validation includes 2016-02-29) \\
Horizon $h$ & 28 days, 8 weekly rolling origins, expanding window \\
M5 sample & 500 series, stratified by archetype $\times$ department \\
Store sample & All 500 series (complete panel) \\
Evaluation population & $500 \times 8 \times 28 = 112{,}000$ forecasts per model--dataset \\
Seed & 42 everywhere \\
\bottomrule
\end{tabular}
\end{table}

\subsection{Series selection}
\label{sec:selection}
Store demand uses all 500 series --- no sampling needed. M5 contributes a
stratified sample of 500 (recorded in
\texttt{05\_experiments/m5\_series\_selection.json}) covering all demand
archetypes and departments, because random sampling could miss rare but
important demand shapes and full-catalog deep learning was computationally
prohibitive. \Cref{fig:representativeness} documents representativeness.

\begin{figure}[H]
\centering
\includegraphics[width=0.9\textwidth]{01_m5_representativeness.png}
\caption{M5 500-series stratified selection: coverage across archetypes and
departments. Source: notebook 04.}
\label{fig:representativeness}
\end{figure}

\subsection{Demand characterization and archetypes}
\label{sec:archetypes-method}
Following the Syntetos--Boylan tradition \citep{syntetos2005}, intermittency is
characterized per series by average demand interval (ADI) and squared
coefficient of variation (CV$^2$), yielding quadrants labeled here as Smooth,
Intermittent, Erratic/Variable, and Lumpy/Highly Intermittent, plus a Moderate
band. Volume terciles add a second axis. The resulting M5 archetype counts are
194 highly intermittent, 199 intermittent, 26 moderate, 80 variable, and 1
smooth --- the last being anecdotal by construction (\cref{sec:smooth-cap}).

\begin{figure}[H]
\centering
\includegraphics[width=0.9\textwidth]{02_intermittency_quadrants.png}
\caption{Intermittency quadrants (ADI vs.\ CV$^2$): the map that motivates both
the archetype analysis and the Croston rung. Source: notebook 03.}
\label{fig:quadrants}
\end{figure}

\subsection{Model-ladder design}
\label{sec:ladder}
Models ascend in complexity while the evaluation stays fixed: baselines (Naive,
Seasonal Naive, Moving Average) set the ``beat this to matter'' bar; SES, DES
(Holt), and additive TES add memory for level, trend, and seasonality; ARIMA
adds autocorrelation modeling and SARIMA seasonal autocorrelation; the
Croston family specializes in intermittent demand; the global LSTM learns
nonlinear patterns across all series of a dataset. Complexity is the
independent variable; value is what is measured. \Cref{fig:ladder} sketches the
ladder.

\begin{figure}[H]
\centering
\includegraphics[width=0.85\textwidth]{11_forecasting_model_ladder.png}
\caption{The model ladder: complexity rises rung by rung; whether value rises
with it is the empirical question (it does not, monotonically). Teaching
figure from notebook 05a.}
\label{fig:ladder}
\end{figure}

\subsection{Forecasting models}
\label{sec:forecast-models}
Naive repeats the last observed day; Seasonal Naive repeats the last seven
observed days (an early implementation that peeked beyond the seasonal window
was caught and corrected to repeat only observed values); Moving Average uses
validation-selected global windows (M5: 14 days; Store: 7 days). SES (level,
$\alpha=0.1$), DES (level $+$ trend, $0.2/0.2$), and additive TES (level $+$
trend $+$ weekly season, $0.3/0.3/0.3$) use transparent implementations with
validation-grid selection. ARIMA(1,1,0) runs on both datasets with a sparse-data
pre-check (series with $>85\%$ zeros fall back to Naive rather than fitting
noise). SARIMA(1,1,0)(0,1,1)$_7$ runs Store-only: seasonal differencing on
mostly-zero series is methodologically dubious, and the Croston family already
covers the M5 specialty. Croston/SBA/TSB (validation alphas 0.1 / 0.5 /
0.1--0.3) run M5-only, since a flat intermittent-demand forecast has no premise
on dense data. The LSTM is one global network per dataset (28-day lookback,
per-series history-only scaling, chronological per-series validation), so it
learns cross-series patterns rather than memorizing individual items. Full
configurations live in \cref{app:configs}.

\subsection{Forecast evaluation metrics}
\label{sec:metrics}
Headline accuracy is mean absolute error (MAE); the primary scale-free metric
is MASE with seasonal period $m=7$ (weekly retail rhythm,
\cref{fig:weekly}), whose denominator is the in-sample one-step seasonal-naive
error on training data only --- validation and test never enter it. RMSE,
RMSSE, WAPE, and sMAPE are reported alongside; no metric is cherry-picked.
Two disciplined choices matter. First, denominators that hit zero (constant
training history) stay honest NaNs instead of epsilon-guarded fakes (zero
occurred nowhere: $n_{\text{NaN}}=0$ everywhere). Second, archetype-level WAPE
and sMAPE are \emph{pooled} ratios over each cell, because averaging per-series
ratios lets one near-zero-demand series explode the mean (an intermediate
computation hit $\sim 2\times10^9$ before the fix); model-level WAPE columns
remain per-series means and were verified consistent across files. sMAPE is
reported with its known zero-demand distortions documented, not hidden.
Definitions are formalized in \cref{app:metrics}.

\begin{figure}[H]
\centering
\includegraphics[width=0.9\textwidth]{03_weekly_seasonality.png}
\caption{Weekly seasonality in retail demand: the business-week rhythm that
justifies seasonal period $m=7$ in MASE/RMSSE, Seasonal Naive, and SARIMA.
Source: notebook 03.}
\label{fig:weekly}
\end{figure}
\begin{figure}[H]
\centering
\includegraphics[width=0.9\textwidth]{04_acf_pacf.png}
\caption{Autocorrelation structure: demand remembers its recent past, which is
the premise ARIMA-style models exploit. Source: notebook 03.}
\label{fig:acf}
\end{figure}

\subsection{Statistical comparison}
\label{sec:stats-method}
Ninety-one paired comparisons (55 M5 $+$ 36 Store) combine the Wilcoxon
signed-rank test (normal approximation) with the Diebold--Mariano test
\citep{diebold1995}, Holm correction \citep{holm1979} for multiplicity,
500-replication bootstrap confidence intervals \citep{efron1993}, Cohen's paired $d_z$, and the
Wilcoxon rank-biserial $r$. With 112,000 paired errors, nearly everything is
``significant'' --- which is exactly why every $p$-value travels with an effect
size. Methodology detail is in \cref{app:stats}.

\subsection{Inventory policy}
\label{sec:policy}
All forecasts drive one shared ordering engine
(\texttt{11\_src/inventory\_policy.py}), identical for every model by
construction: a daily-review order-up-to (base-stock) policy with lost sales.
Per series--origin instance, safety stock is $z$ times the forecast-error
standard deviation (floored) times $\sqrt{L}$; the order-up-to level is the sum
of the next $L$ forecasts plus safety stock; unmet demand is lost at penalty
$P$ per unit; end-of-day stock accrues holding cost $H$ per unit. The default
ranking policy is $L=7$, 95\% service target ($z\approx1.645$), $H=1$, $P=5$.
Costs decompose into holding $+$ stockout, with service level, average
inventory, and reorder counts as diagnostics. \Cref{fig:invlink} shows how
forecasts become orders.

\begin{figure}[H]
\centering
\includegraphics[width=0.85\textwidth]{12_inventory_connection.png}
\caption{From forecast to shelf: forecasts feed safety stock and order-up-to
levels; the policy --- not the error metric --- decides what an error costs.
Teaching figure from notebook 05a.}
\label{fig:invlink}
\end{figure}

\subsection{Robustness and sensitivity analysis}
\label{sec:sens-method}
Headline rankings use the default policy, but trust requires more: a 27-policy
grid ($L \in \{3,7,14\} \times$ service $\in \{90\%,95\%,99\%\} \times P \in
\{3,5,10\}$, $H=1$) replays every model--dataset through 540 costing runs. A
model first across nearly all policies is \emph{robust}; one whose crown moves
with the policy is \emph{fragile}. Per-archetype M5 costing adds the
demand-structure dimension.

\subsection{Convergence and feasibility}
\label{sec:conv-method}
Every ARIMA/SARIMA fit is logged per series--origin
(\texttt{convergence\_details\_*.csv}) with a scope column separating the
exploratory 100-series SARIMA subset from the primary full-500 run. A
16-fit feasibility probe (0.62\,s/fit) preceded the 4,000-fit full-500 SARIMA
run (2,382\,s, 4,000/4,000 successful). Zero-failure convergence on dense data
is itself a finding when contrasted with the sparse-data fallback rate.

\subsection{Validation and reproducibility}
\label{sec:validation}
An independent validator re-checks row counts (112,000 per model--dataset),
duplicate keys (zero), origin coverage (1--8), horizon length (28 everywhere),
test-window containment, forecast finiteness/non-negativity, and cross-family
actuals identity (max spread 0.0): 0 FAIL
(\texttt{06\_results/validation\_report.json}). Seeds (42), manifests, run
logs, and a pre-hardening backup make every step replayable; provenance is
inventoried in \cref{app:repro}.

\subsection{Research hardening and final evidence freeze}
\label{sec:hardening}
Before reporting, the evidence passed a hardening audit that corrected the
calendar (leap day 2016-02-29 assigned to validation: 1,034/121/83),
revalidated the LSTM on chronological per-series splits, fixed a Seasonal Naive
lookahead, unified all models onto one inventory module, corrected MASE
grouping and Wilcoxon effect sizes, pooled archetype WAPE/sMAPE, replaced the
SARIMA subset with the full-500 run as the primary result, and repaired
identifiers. Pre-hardening originals are preserved untouched in
\texttt{06\_results/\_pre\_hardening\_backup/} and never used as report
sources. The freeze (commits 53c8343 through edb42ad) generated the number
sheet, figure/table audits, and claim traceability programmatically --- no
headline number in this report was typed by hand. Hardening is reported to
explain why the evidence can be trusted, not as a debugging diary.
